\documentclass{article}

\usepackage[utf8]{inputenc}
\usepackage[turkish, shorthands = :!]{babel}
\usepackage[T1]{fontenc}
\usepackage{lmodern}

\usepackage{fancyhdr, lipsum}

\usepackage{graphicx} % Enhanced support for graphics (resim ekleyebilmek için)

\usepackage{tabularx} % metnin yanına resim koyma yollarından biri

\graphicspath{{./auxiliary/}} % resimler auxiliary klasörü içinde

\renewcommand{\headrulewidth}{0 pt}
\pagestyle{fancy}
\fancyhf{}
\fancyhead[C]{\thepage}

\begin{document}

%%%%%%% SAYFA 1 %%%%%%%

\lipsum[1][1-2]
\includegraphics[scale = 0.1]{logo}
\lipsum[1][3-4]

\bigskip
\hrule
\bigskip

\lipsum[1][1-2]

\includegraphics[scale = 0.1]{logo}

\lipsum[1][3-4]

\bigskip
\hrule
\bigskip

\lipsum[1][1-2]

\begin{center}
\includegraphics[scale = 0.1]{logo}
\end{center}

\lipsum[1][3-4]

\newpage

%%%%%%% SAYFA 2 %%%%%%%

\lipsum[1][1-2]
% burada boşluk olup olmaması, dikey boşlukları değilse de, resimden önceki ve sonraki metni etkileyebiliyor
\begin{figure}[h] % 'h'ere (özellikle gerekmiyorsa kullanmamak daha iyi olabilir)
\centering
\includegraphics[scale = 0.1]{logo}
\caption{Medeniyet Üniversitesi logosu.} % figürlerde başlık aşağıda, tablolarda yukarıda; bk. https://tex.stackexchange.com/questions/3243/why-should-a-table-caption-be-placed-above-the-table
\label{fig:logo}
\end{figure}
% burada boşluk olup olmaması bir şey değiştirmiyor
\lipsum[1][3-4]

Şekil~\ref{fig:logo}...

\bigskip
\hrule
\bigskip

\lipsum[1][1-2]

\begin{figure}[h]
\begin{center} % \centering yerine bunu kullanmak fazladan dikey boşluk yaratıyor (yukarıdaki örneklerde de gördüğümüz gibi)
\includegraphics[scale = 0.1]{logo}
\end{center}
\caption{Medeniyet Üniversitesi logosu.}
\end{figure}

\lipsum[1][3-4]

\newpage

%%%%%%% SAYFA 3 %%%%%%%

\lipsum[1][1-5]

\bigskip
\hrule
\bigskip

\begin{minipage}{\textwidth}
\lipsum[1][1-5]
\end{minipage}

\bigskip
\hrule
\bigskip

\noindent
\begin{minipage}{\textwidth}
\lipsum[1][1-5]
\end{minipage}

\bigskip
\hrule
\bigskip

\noindent
\fbox
{
\begin{minipage}{\textwidth}
\lipsum[1][1-5]
\end{minipage}
}

\bigskip
\hrule
\bigskip

\noindent
\frame % \frame'de \fboxsep söz konusu değil
{
\begin{minipage}{\textwidth}
\lipsum[1][1-5]
\end{minipage}
}

\bigskip
\hrule
\bigskip

\noindent
\frame
{
\begin{minipage}{\textwidth}
\lipsum[1][1-5]
\end{minipage}% bu yüzde işareti (ya da müteakip askı parantez) tam burada olmalı, yoksa çerçeve \textwidth'i az da olsa aşıyor 
}

\newpage

%%%%%%% SAYFA 4 %%%%%%%

\lipsum[1][1-5]

\bigskip
\hrule
\bigskip

\noindent
\frame
{
\begin{minipage}{0.5\textwidth}
\lipsum[1][1-2]
\end{minipage}
}
\frame
{
\begin{minipage}{0.5\textwidth}
\lipsum[1][1-2]
\end{minipage}
}

\bigskip
\hrule
\bigskip

% \textwidth'in aşılmaması için aşağıdaki üç yüzde işareti de elzem
\noindent
\frame
{
\begin{minipage}{0.5\textwidth}
\lipsum[1][1-2]
\end{minipage}%
}%
\frame
{
\begin{minipage}{0.5\textwidth}
\lipsum[1][1-2]
\end{minipage}%
}

\bigskip
\hrule
\bigskip

\noindent
\frame
{
\begin{minipage}{0.5\textwidth}
\lipsum[1][1-4]
\end{minipage}%
}%
\frame
{
\begin{minipage}{0.5\textwidth}
\centering
\includegraphics[scale = 0.1]{logo}
\end{minipage}%
}

\newpage

%%%%%%% SAYFA 5 %%%%%%%

\lipsum[1][1-5]

\bigskip
\hrule
\bigskip

% resim metnin üst kısmıyla aynı hizada değil (metin çok kısa olsa da sıkıntılı)
\noindent
\begin{minipage}{0.5\textwidth}
\lipsum[1][1-8]
\end{minipage}
\hfill
\includegraphics[scale = 0.1]{logo}

\bigskip
\hrule
\bigskip

% resmin metnin üst kısmıyla aynı hizada olmasını sağlayan ilk çözüm (bk. 2018 Güz MAT 101 finali); \textwidth'in aşılmaması için aşağıdaki yüzde işaretleri elzem
\noindent
\begin{minipage}[t]{0.5\textwidth}
\vspace{0 cm}
\lipsum[1][1]
\end{minipage}%
%
\begin{minipage}[t]{0.5\textwidth}
\vspace{0 cm}
\begin{minipage}{\textwidth}
\begin{flushright}
\includegraphics[scale = 0.1]{logo}
\end{flushright}
\end{minipage}
\end{minipage}

\bigskip
\hrule
\bigskip

% bir başka çözüm (fazladan dikey boşluk \vspace{0 cm} satırları yüzünden ama onlarsız da olmuyor); krş. https://tex.stackexchange.com/questions/11630/aligning-image-and-text-on-top-with-minipages
\begingroup
\setlength{\tabcolsep}{0 pt} % yoksa biraz içerlek duruyor metin
\noindent
\begin{tabularx}{\textwidth}{
	>{\hsize = 1\hsize }X
	>{\hsize = 1\hsize \raggedleft\arraybackslash}X
}
    \vspace{0 pt} 
    \lipsum[1][1]
    & 
    \vspace{0 pt}
    \includegraphics[scale = 0.1]{logo}
\end{tabularx}
\endgroup

\newpage

%%%%%%% SAYFA 6 %%%%%%%

\lipsum[1][1-5]

\bigskip
\hrule
\bigskip

\raisebox{0.5 cm}{\includegraphics[scale = 0.1]{logo}}
\includegraphics[scale = 0.1]{logo}

\bigskip
\hrule
\bigskip

\begin{enumerate}

\item \lipsum[1][1-3]

\item
% burada \noindent'e gerek olmuyor
% aşağıda \linewidth yerine \textwidth dersek dışarı taşıyoruz
\begin{minipage}[t]{0.5\linewidth}
\vspace{0 cm}
\lipsum[1][1]
\end{minipage}%
%
\begin{minipage}[t]{0.5\linewidth}
\vspace{0 cm}
\begin{minipage}{\linewidth}
\begin{flushright}
\includegraphics[scale = 0.1]{logo}
\end{flushright}
\end{minipage}
\end{minipage}

% krş. https://tex.stackexchange.com/a/30369/178454
\item
\raisebox{0.56\baselineskip} % büyük bir harfin boyunun (mesela L'nin uzun kenarının) \baselineskip'e (yani ardışık satırların "base line"ları arasındaki uzaklığa) oranını ekrandan cetvelle ölçerek yaklaşık 0,56 bulduk
{% bu yüzde işareti elzem
\begin{minipage}[t]{0.5\linewidth}
\vspace{0 cm}
\lipsum[1][1]
\end{minipage}%
%
\begin{minipage}[t]{0.5\linewidth}
\vspace{0 cm}
\begin{minipage}{\linewidth}
\begin{flushright}
\includegraphics[scale = 0.1]{logo}
\end{flushright}
\end{minipage}
\end{minipage}% bu yüzde işareti elzem
}

\item \lipsum[1][1-3]

\end{enumerate}

\end{document}